%% LyX 2.4.4 created this file.  For more info, see https://www.lyx.org/.
%% Do not edit unless you really know what you are doing.
\documentclass[english]{article}
\usepackage[T1]{fontenc}
\usepackage[latin9]{inputenc}
\usepackage{enumitem}
\usepackage{amsmath}
\usepackage{amsthm}
\usepackage{geometry}
\geometry{verbose,tmargin=1in,bmargin=1in,lmargin=1in,rmargin=1in}

\makeatletter
%%%%%%%%%%%%%%%%%%%%%%%%%%%%%% Textclass specific LaTeX commands.
\numberwithin{equation}{section}
\numberwithin{figure}{section}
\newlength{\lyxlabelwidth}      % auxiliary length 

%%%%%%%%%%%%%%%%%%%%%%%%%%%%%% User specified LaTeX commands.
\usepackage{fancyhdr}
\usepackage{lastpage}
\pagestyle{fancy}
\usepackage{enumitem}
\usepackage{ifthen}
%sets math fonts to be more readable
\everymath=\expandafter{\the\everymath\displaystyle}
%\DeclareMathSizes{10}{10}{10}{10}
\usepackage{multicol}
% Added by lyx2lyx
\setlength{\parskip}{\smallskipamount}
\setlength{\parindent}{0pt}

\makeatother

\usepackage{babel}
\begin{document}
\renewcommand{\author}{Dr.\,James Ashe}

\newcommand{\classNum}{232}

\newcommand{\classSec}{C}

\newcommand{\nameType}{}

\newcommand{\examType}{Exam 3 Review}

\renewcommand{\date}{Exam 3 will be be given on December $5^{\text{th}}$, 2012} %%\currenttime, \today

%%First page's header%%%%%%%%%%%%%%%%%%%%%%
\fancypagestyle{firststyle} %%header settings for first pagr
{
	\fancyhead[L]{MTH \classNum{}(\classSec{}) \author{}\\
	\textbf{\examType{}} \date}
	\fancyhead[C]{\textbf{\nameType}}
	\setlength{\headsep}{14pt}
}
%%%%%%%%%%%%%%%%%%%%%%%%%%%%%%%%%%%%%%%%%%

%%Next page's header%%%%%%%%%%%%%%%%%%%%%%%
\setlist{nolistsep}
\lhead{\classNum{}(\classSec{})} 
\chead{\textbf{\nameType}} 
\cfoot{\thepage{}~of~\pageref{LastPage}} 
\setlength{\headsep}{5pt} 
%%%%%%%%%%%%%%%%%%%%%%%%%%%%%%%%%%%%%%%%%%%

%%Various settings; see the preamble for other settings%%%%%
%%\setenumerate[0]{leftmargin=12pt,itemindent=0pt,labelwidth=0pt}
\renewcommand{\labelenumi}{\textbf{\arabic{enumi}.}}
\renewcommand{\labelenumii}{\textbf{\alph{enumii}.}}
\thispagestyle{firststyle}
%%%%%%%%%%%%%%%%%%%%%%%%%%%%%%%%%%%%%%%%%%

\vfill{}

\subsection*{7.1: Inverse Functions}
\begin{itemize}
\item Find the inverse, $f^{-1}(x)$, of a function on a given interval
and its domain. 

\begin{itemize}
\item Directly; \textbf{Problems 15-19 and 25-27.}
\end{itemize}
\item Verify the relationship, i.e., \textbf{$f(f^{-1}(x))=x$ }and $f^{-1}(f(x))=x$;
\textbf{Problems 15-19.}
\item Graph the inverse function.

\begin{itemize}
\item From a function; \textbf{Problems 25-27.}
\item From a graph; \textbf{Problems 29-30 (these were not assigned).}
\end{itemize}
\item Find the derivative of an inverse function using $\left(f^{-1}\right)(y_{0})=\frac{1}{f'(x_{0})}$
where $y_{0}=x_{0}$; \textbf{Problems 35-37.}\vfill
\end{itemize}

\subsection*{7.2: Differentiation and Integrating the Logarithmic and Exponential
Functions}
\begin{itemize}
\item Find the derivatives of functions with $\ln x$, $\ln\left(u(x)\right)$,
$e^{x}$, and $e^{u(x)}$;\textbf{ Problems 15-19 and 25-27.}

\begin{itemize}
\item You may have to use the product rule, quotient rule, and/or chain
rule.
\end{itemize}
\item Find the indefinite and definite integrals of functions with $e^{x}$,
$e^{u(x)}$, $\frac{1}{x}$, and $\frac{1}{u(x)}$ (not $\ln x$);
\textbf{Problems 31-35, 37.}

\begin{itemize}
\item You may have to use substitution. \vfill
\end{itemize}
\item You will NOT be given problems where you have to use logarithmic differentiation,
$\frac{d}{dx}\ln f(x)=f'(x)\frac{1}{f(x)}$ (Problems 39-41).
\end{itemize}

\subsection*{7.3: Logarithmic and Exponential Functions with other bases. }
\begin{itemize}
\item Understanding log functions; \textbf{Problems 11-15.}
\item Finding the derivative and indefinite integrals of logarithmic and
exponential functions with bases other than $e$\textbf{.}

\begin{itemize}
\item Derivatives of $b^{x}$; \textbf{Problems 21-25.}
\item Indefinite Integrals with general bases; \textbf{Problems 30-33 (indefinite
only).}
\end{itemize}
\item Find a derivative using the general power rule. \textbf{Problems 35,
36, 42.}

\begin{itemize}
\item The results of section 7.3 allows the power rule (early stated for
rational powers) to be extended to include all real numbers. 
\end{itemize}
\item Derivative and integral Exam problems from this section will be straight
forward application of the formulas. There will probably be at most
one problem taken from each item above.\vfill
\end{itemize}

\subsection*{8.1 Integration by Parts.}
\begin{itemize}
\item Find the indefinite integral of a function using integration by parts;
\textbf{Problems 7-19.}

\begin{itemize}
\item You may have to use substitution. You will NOT have to use repeated
integration by parts (Problems 23-28).
\end{itemize}
\item Find the definite integral of a function using integration by parts;
\textbf{Problems 29-33.}

\begin{itemize}
\item Some of the HW problems may require trigonometric substitutions found
in earlier problems (e.g., 35). You will NOT have to use trigonometric
substitutions on the exam; we will cover this in 8.2 and 8.3.\vfill
\end{itemize}
\end{itemize}

\subsection*{Some general notes.}
\begin{itemize}
\item You will be given the necessary formulas from section 7.1, 7.2, 7.3,
and 8.1 (as they appear in the text). You will also be given the product,
quotient, chain, and trigonometric rules.
\item Of the four sections, 7.3 will be covered the least, however be prepared
to use the formulas from that section.\vfill
\item There will be 15-20 problems including some ``easy'' problems involving
properties of logs and inverse functions.
\end{itemize}

\end{document}
